% Part: first-order-logic
% Chapter: introduction
% Section: models-theories

\documentclass[../../../include/open-logic-section]{subfiles}

\begin{document}

\olfileid{fol}{int}{mod}

\olsection{Modellen en theorieën}

Zodra we de syntaxis en semantiek van de eerste-ordelogica hebben
gedefinieerd, kunnen we de eigenschappen van !!{structure}s en de
semantische begrippen onderzoeken. We kunnen ook
!!{derivation}ssystemen definiëren en bestuderen. Bij een verzameling
!!{sentence}s kunnen we vragen welke !!{structure}s alle !!{sentence}s
uit die verzameling waar maken. Gegeven een verzameling
!!{sentence}s~$\Gamma$ heet !!a{structure}~$\Struct{M}$ die ze vervult
een \emph{model van~$\Gamma$}. We kunnen uitgaan van~$\Gamma$ en haar
modellen proberen te vinden: hoe zien ze eruit en hoe groot of klein
moeten ze zijn? We kunnen daarentegen ook uitgaan van één !!{structure}
of een klasse van !!{structure}s en vragen welke !!{sentence}s daarin
waar zijn. Bestaan er !!{sentence}s die deze !!{structure}s
\emph{karakteriseren}, in die zin dat zij daarin en uitsluitend daarin
waar zijn? Zulke vragen behoren tot de \emph{modeltheorie}. Zij liggen
ook ten grondslag aan de \emph{axiomatische methode}: een klasse van
!!{structure}s beschrijven door middel van een verzameling
!!{sentence}s, de axioma's van een theorie. Dit is mogelijk doordat
precies de !!{sentence}s die in de eerste-ordelogica semantisch uit de
axioma's volgen, waar zijn in alle modellen van die axioma's.

Beschouw als zeer eenvoudig voorbeeld preorden. Een preorde is een
relatie~$R$ op een verzameling~$A$ die zowel reflexief als transitief
is. Een verzameling~$A$ met daarop een tweeplaatsige relatie
$R \subseteq A \times A$ is precies wat nodig is voor !!a{structure}
voor een eerste-ordetaal met één tweeplaatsig relatiesymbool~$\Obj P$:
we stellen $\Domain{M} = A$ en $\Assign{\Obj P}{M} = R$. Omdat $R$ een
preorde is, is zij reflexief en transitief. We kunnen een
verzameling~$\Gamma$ van !!{sentence}s van de eerste-ordelogica geven
die dit uitdrukken:
\begin{align*}
  & \lforall[\Obj v_0][\Atom{\Obj P}{\Obj v_0,\Obj v_0}]\\
  & \lforall[\Obj v_0][\lforall[\Obj v_1][\lforall[\Obj v_2][((\Atom{\Obj P}{\Obj v_0,\Obj v_1} \land \Atom{\Obj P}{\Obj v_1,\Obj v_2}) \lif \Atom{\Obj P}{\Obj v_0,\Obj v_2})]]]
\end{align*}
Deze !!{sentence}s zijn eenvoudigweg de symboliseringen van ``voor
iedere~$x$ geldt $Rxx$'' ($R$ is reflexief) en ``als $Rxy$ en $Ryz$,
dan ook $Rxz$'' ($R$ is transitief). We zien dat
!!a{structure}~$\Struct{M}$ een model van deze twee
!!{sentence}s~$\Gamma$ is dan en slechts dan als $R$ (dus
$\Assign{\Obj P}{M}$) een preorde op~$A$ (dus $\Domain{M}$) is. De
modellen van~$\Gamma$ zijn met andere woorden precies de preorden.
Iedere eigenschap van alle preorden die kan worden uitgedrukt in de
eerste-ordetaal met alleen~$\Obj P$ als !!{predicate} (zoals de
reflexiviteit en transitiviteit hierboven), volgt semantisch uit de
twee !!{sentence}s in~$\Gamma$, en omgekeerd. Wat we over de modellen
van~$\Gamma$ bewijzen, hebben we dus over alle preorden bewezen.

Bij iedere specifieke theorie en klasse van modellen, zoals $\Gamma$
en alle preorden, rijzen interessante vragen over wat wel en niet in de
bijbehorende eerste-ordetaal kan worden uitgedrukt. Sommige
eigenschappen van !!{structure}s zijn voor alle talen en klassen van
modellen van belang, namelijk de eigenschappen die de grootte van het
!!{domain} betreffen. Men kan bijvoorbeeld uitdrukken dat het
!!{domain} precies $n$~!!{element}s bevat, voor iedere~$n \in \PosInt$.
Met een
oneindige verzameling !!{sentence}s kan men ook uitdrukken dat het
!!{domain} oneindig is. Men kan echter niet uitdrukken dat het domein
eindig of !!{nonenumerable} is. Deze resultaten over de beperkingen van
eerste-ordetalen zijn gevolgen van de compactheidsstelling en de
stelling van L\"owenheim--Skolem.

\end{document}
